% !TeX root = ../main.tex
\begin{frame}{课程评价——刘诗扬}
  \begin{SlideCanvas}
    \node[anchor=north west,inner sep=0pt,text width=850pt,align=left,
      text=BICNavy] at (55,80) {\begin{minipage}{850pt}
      \usebeamerfont{BIC detail}\raggedright
      \setlength{\parindent}{0pt}\setlength{\parskip}{4pt}
      \setlength{\columnsep}{26pt}\setlength{\multicolsep}{0pt}
      \clubpenalty=10000\widowpenalty=10000
\begin{multicols}{2}\raggedcolumns
\par\addvspace{6pt}{\usebeamerfont{BIC point}一、课程评价\par}\nobreak
本课程内容丰富，涵盖医学影像、人工智能、医疗机器人、生物医学材料及MEMS传感器等多个前沿领域，充分体现了医工交叉学科的特点。授课教师结合自身科研方向，介绍了相关技术的基本原理、研究进展及实际应用，使课程兼具理论深度与实践价值。整体课程安排合理，内容紧跟科技发展趋势，有效拓宽了学生的学术视野。\par
\par\addvspace{6pt}{\usebeamerfont{BIC point}二、学习收获\par}\nobreak
通过本课程的学习，我对医工交叉领域的发展现状及未来趋势有了更加全面的认识，了解了先进材料、智能传感、人工智能及医疗机器人等技术在医学中的应用。同时，进一步认识到材料科学、微电子技术与医学工程之间的紧密联系，以及多学科协同创新的重要性。课程中的科研案例也为我提供了新的研究思路，使我更加重视基础研究与实际应用需求的结合，对今后的科研工作具有积极的指导意义。\par
\par\addvspace{6pt}{\usebeamerfont{BIC point}三、意见与建议\par}\nobreak
建议今后的课程在保持前沿性和学科交叉特色的基础上，适当增加随堂讨论与实例展示环节。在随堂讨论方面，可以围绕关键技术、研究难点及未来发展方向设置开放性问题，鼓励学生结合自身专业背景交流观点，激发创新思维。在实例展示方面，建议更多结合教师的实际科研项目和典型医疗应用案例，使理论知识更加直观易懂。同时，适当介绍相关技术从基础研究到临床应用的转化过程，进一步提升课程的参与度和教学效果。\par
\end{multicols}
    \end{minipage}};
  \end{SlideCanvas}
  \note{原始材料为用户提供的四份课程反馈之一，正文逐字保留，仅调整标题、段落和编号排版。刘诗扬《医工交叉研究前沿进展》课程反馈与建议。一、课程评价\par 本课程内容丰富，涵盖医学影像、人工智能、医疗机器人、生物医学材料及MEMS传感器等多个前沿领域，充分体现了医工交叉学科的特点。授课教师结合自身科研方向，介绍了相关技术的基本原理、研究进展及实际应用，使课程兼具理论深度与实践价值。整体课程安排合理，内容紧跟科技发展趋势，有效拓宽了学生的学术视野。\par \par 二、学习收获\par 通过本课程的学习，我对医工交叉领域的发展现状及未来趋势有了更加全面的认识，了解了先进材料、智能传感、人工智能及医疗机器人等技术在医学中的应用。同时，进一步认识到材料科学、微电子技术与医学工程之间的紧密联系，以及多学科协同创新的重要性。课程中的科研案例也为我提供了新的研究思路，使我更加重视基础研究与实际应用需求的结合，对今后的科研工作具有积极的指导意义。\par \par 三、意见与建议\par 建议今后的课程在保持前沿性和学科交叉特色的基础上，适当增加随堂讨论与实例展示环节。在随堂讨论方面，可以围绕关键技术、研究难点及未来发展方向设置开放性问题，鼓励学生结合自身专业背景交流观点，激发创新思维。在实例展示方面，建议更多结合教师的实际科研项目和典型医疗应用案例，使理论知识更加直观易懂。同时，适当介绍相关技术从基础研究到临床应用的转化过程，进一步提升课程的参与度和教学效果。}
\end{frame}
